\documentclass[12pt,a4paper]{article}
\usepackage{expediente}
\begin{document}
% =============================================================================
% FICHA DE DATOS DEL EXPEDIENTE --- actualizado 20 ago 2026
% =============================================================================

% --- VENDEDOR / TITULAR REGISTRAL -------------------------------------------
\newcommand{\VendedorNombres}{RICHAR DARWIN}
\newcommand{\VendedorApellidos}{CUTIRE CONDORI}
\newcommand{\VendedorNombreCompleto}{RICHAR DARWIN CUTIRE CONDORI}
\newcommand{\VendedorDNI}{42610690}
\newcommand{\VendedorEstadoCivil}{soltero, sin unión de hecho vigente}
\newcommand{\VendedorOcupacion}{\faltante{[ocupación de Richar]}}
\newcommand{\VendedorDomicilio}{\faltante{[domicilio actual de Richar]}}
\newcommand{\VendedorTelefono}{\faltante{[teléfono de Richar]}}
\newcommand{\VendedorCorreo}{\faltante{[correo]}}
\newcommand{\VendedorNacionalidad}{peruano}
\newcommand{\VendedorFechaNac}{\faltante{[fecha de nacimiento]}}
\newcommand{\VendedorDNIEmision}{\faltante{[fecha de emisión del DNI de Richar --- la pide SUNARP]}}

% --- CONVIVIENTE: NO CORRESPONDE -------------------------------------------
\newcommand{\ConvivienteNombre}{NO CORRESPONDE}
\newcommand{\ConvivienteDNI}{NO CORRESPONDE}
\newcommand{\ConvivienteDesde}{NO CORRESPONDE}
\newcommand{\ConvivienteInscrita}{no existe unión de hecho vigente}
\newcommand{\ConvivienteDomicilio}{NO CORRESPONDE}

% --- COMPRADORES (copropiedad en partes iguales) ---------------------------
\newcommand{\CompradorUnoNombre}{NICO ALVARO CUTIRE ARCE}
\newcommand{\CompradorUnoDNI}{48325017}
\newcommand{\CompradorUnoEmision}{01 de enero de 2024}
\newcommand{\CompradorUnoEstadoCivil}{soltero}
\newcommand{\CompradorUnoOcupacion}{ingeniero informático}
\newcommand{\CompradorUnoDomicilio}{Jr.\ Wiracocha, Sicuani, provincia de Canchis, Cusco, Perú}
\newcommand{\CompradorUnoTelefono}{984715108}

\newcommand{\CompradorDosNombre}{LUCY CUTIRE ARCE}
\newcommand{\CompradorDosDNI}{47478852}
\newcommand{\CompradorDosEmision}{09 de enero de 2026}
\newcommand{\CompradorDosEstadoCivil}{soltera}
\newcommand{\CompradorDosOcupacion}{contador}
\newcommand{\CompradorDosDomicilio}{Comunidad Hanocca, distrito de Layo, provincia de Canas, Cusco, Perú}
\newcommand{\CompradorDosTelefono}{953440875}

\newcommand{\CompradorNombres}{NICO ALVARO y LUCY}
\newcommand{\CompradorApellidos}{CUTIRE ARCE}
\newcommand{\CompradorNombreCompleto}{NICO ALVARO CUTIRE ARCE y LUCY CUTIRE ARCE}
\newcommand{\CompradorDNI}{48325017 y 47478852}
\newcommand{\CompradorEstadoCivil}{solteros}
\newcommand{\CompradorConyuge}{NO CORRESPONDE}
\newcommand{\CompradorConyugeDNI}{NO CORRESPONDE}
\newcommand{\CompradorOcupacion}{ingeniero informático y contador, respectivamente}
\newcommand{\CompradorDomicilio}{Sicuani (Canchis) y Layo (Canas), Cusco}
\newcommand{\CompradorTelefono}{984715108 / 953440875}
\newcommand{\CompradorCorreo}{\faltante{[correo de contacto]}}
\newcommand{\CompradorNacionalidad}{peruanos}
\newcommand{\PorcentajeCopropiedad}{50\,\% cada uno}

% --- INTERMEDIARIO / OCUPANTE DEL REMANENTE --------------------------------
\newcommand{\IntermediarioNombre}{LUCIO GIL CUTIRE PARI}
\newcommand{\IntermediarioDNI}{42467850}
\newcommand{\IntermediarioEstadoCivil}{\faltante{[estado civil de Lucio]}}
\newcommand{\IntermediarioOcupacion}{\faltante{[ocupación]}}
\newcommand{\IntermediarioDomicilio}{\faltante{[domicilio de Lucio]}}
\newcommand{\IntermediarioTelefono}{\faltante{[teléfono de Lucio]}}
\newcommand{\IntermediarioParentesco}{familiar (apellido Cutire)}

% --- PREDIO MATRIZ ----------------------------------------------------------
\newcommand{\PartidaMatriz}{\faltante{[N.º de partida electrónica SUNARP --- pendiente]}}
\newcommand{\OficinaRegistral}{Oficina Registral de Puerto Maldonado}
\newcommand{\AreaMatriz}{\faltante{[área total inscrita; las partes refieren del orden de 100 ha]}}
\newcommand{\UnidadCatastral}{\faltante{[código catastral / UC, si existe]}}
\newcommand{\NombreFundo}{\faltante{[nombre del fundo, si tiene]}}
\newcommand{\KmCarretera}{Interoceánica Sur, al este del Puente Jayave (el puente queda al oeste del predio, sobre el río; no debajo del lote)}

% --- PREDIO A INDEPENDIZAR --------------------------------------------------
\newcommand{\AreaIndependizarHa}{10{,}00}
\newcommand{\AreaIndependizarMetros}{100 000{,}00}
\newcommand{\PerimetroM}{2 200{,}58}
\newcommand{\FrenteM}{100{,}29}
\newcommand{\FondoM}{1 000{,}00}

\newcommand{\PUnoLat}{$-12{,}912452$}
\newcommand{\PUnoLon}{$-70{,}170058$}
\newcommand{\PDosLat}{$-12{,}912497$}
\newcommand{\PDosLon}{$-70{,}170981$}
\newcommand{\PTresLat}{$-12{,}903469$}
\newcommand{\PTresLon}{$-70{,}171438$}
\newcommand{\PCuatroLat}{$-12{,}903424$}
\newcommand{\PCuatroLon}{$-70{,}170515$}

\newcommand{\PUnoE}{373 065{,}00}
\newcommand{\PUnoN}{8 572 256{,}00}
\newcommand{\PDosE}{372 965{,}00}
\newcommand{\PDosN}{8 572 251{,}00}
\newcommand{\PTresE}{372 911{,}00}
\newcommand{\PTresN}{8 573 249{,}00}
\newcommand{\PCuatroE}{373 011{,}00}
\newcommand{\PCuatroN}{8 573 255{,}00}

\newcommand{\FechaPago}{15 de agosto de 2018}
\newcommand{\MontoPago}{S/\ 20 000{,}00}
\newcommand{\MontoPagoLetras}{VEINTE MIL Y 00/100 SOLES}
\newcommand{\FormaPagoHistorica}{dinero en efectivo, sin medio de pago del sistema financiero}

% Nombres genéricos --- sustituir por los reales al ir a notaría
\newcommand{\NotarioNombre}{Dr.\ JUAN CARLOS RÍOS SALINAS (nombre genérico a sustituir)}
\newcommand{\NotariaCiudad}{Puerto Maldonado}
\newcommand{\ColegiaturaAbogado}{CAL Madre de Dios N.º 0000 (genérico)}
\newcommand{\AbogadoNombre}{Abog.\ ANA MARÍA QUISPE HUAMÁN (nombre genérico a sustituir)}
\newcommand{\IngenieroNombre}{Ing.\ PEDRO ANTONIO MAMANI CONDORI (nombre genérico a sustituir)}
\newcommand{\IngenieroCIP}{CIP 000000 (genérico)}
\newcommand{\Municipalidad}{Municipalidad Distrital de Inambari}
\newcommand{\Provincia}{Tambopata}
\newcommand{\Departamento}{Madre de Dios}
\newcommand{\Distrito}{Inambari}

\newcommand{\LugarOtorgamiento}{Puerto Maldonado}
\newcommand{\DiaOtorgamiento}{\faltante{[día]}}
\newcommand{\MesOtorgamiento}{\faltante{[mes]}}
\newcommand{\AnioOtorgamiento}{2026}

\newcommand{\TestigoUno}{\faltante{[testigo 1: nombres, DNI, domicilio]}}
\newcommand{\TestigoDos}{\faltante{[testigo 2: nombres, DNI, domicilio]}}

\InicioDocumento{Qué son el predial, la alcabala, la notaría y SUNARP en \emph{este} caso}{DOC-23}

\noindent No son trámites abstractos. En esta compraventa de 10 ha en Inambari, cada uno tiene un rol concreto.

\Clausula{1. Impuesto predial (autoavalúo) atrasado}
Es el impuesto anual municipal sobre el terreno. Lo debe quien figura como dueño al 1 de enero de cada año. Como Richar sigue siendo el titular registral desde 2018 hasta hoy, \textbf{la Municipalidad de Inambari puede cobrarle 2018--2026} (y recargos), aunque ustedes ya le hayan pagado las 10 ha en efectivo.

\textbf{En este caso:} el notario suele pedir una constancia de que no hay deuda, o el pago, antes de firmar. Si Richar nunca declaró el predio, hay que empadronarlo y liquidar. Pacten por escrito quién pone el dinero: lo justo es que lo pague \textbf{Richar} (él fue el dueño de registro); si él no puede, Nico y Lucy pueden pagarlo para destrabar, y descontarlo o dejar constancia de adelanto.

No es un impuesto ``de la compraventa''. Es la deuda vieja del fundo.

\Clausula{2. Alcabala}
Es el impuesto municipal \textbf{de la transferencia} (3\,\% sobre el mayor valor entre el precio y el autoavalúo, descontadas las primeras 10 UIT del año). Lo paga el \textbf{comprador}. En 2026 esas 10 UIT equivalen a S/\ 55\,000.

\textbf{En este caso:} el precio es S/\ 20\,000. Si el autoavalúo de las 10 ha no supera S/\ 55\,000, es probable que la liquidación salga \textbf{inafecta o en cero} (el tramo de 10 UIT cubre la base). Igual hay que \textbf{liquidar y sacar constancia}: sin ese papel el notario y SUNARP no inscriben. Si el municipio valúa las 10 ha por encima de S/\ 55\,000, Nico y Lucy pagarían el 3\,\% del exceso.

No se pida liquidar como donación.

\Clausula{3. Notaría}
El notario de \textbf{Puerto Maldonado} identifica a las partes, lee la escritura, toma huellas, deja constancia de que el precio de 2018 \textbf{no se bancarizó}, protocoliza planos y minutas, y envía el parte por SID-SUNARP.

\textbf{En este caso:} se paga un arancel notarial (independización + compraventa, testimonios, copias). El nombre \NotarioNombre\ es \textbf{genérico}: eligen ustedes la notaría real. Sin notario \textbf{no hay inscripción} de las 10 ha a nombre de Nico y Lucy.

\Clausula{4. SUNARP (derechos registrales)}
SUNARP cobra por calificar e inscribir: abrir partida de las 10 ha y registrar la venta a los dos compradores.

\textbf{En este caso:} sin esa inscripción, Richar sigue apareciendo como dueño frente a bancos, herederos y terceros. El papel de 2018 no basta. Los derechos se calculan sobre el valor del acto (precio / autoavalúo). Es un pago distinto de la alcabala.

\Clausula{Quién suele pagar, en esta familia}
{\small\renewcommand{\arraystretch}{1.25}
\noindent\begin{tabularx}{\textwidth}{>{\bfseries}p{4.2cm} X}
Predial 2018--2026 & Richar (titular). Nico/Lucy solo si hay que adelantarlo para firmar. \\
Alcabala 2026 & Nico y Lucy (adquirentes). Posible inafectación si la base $<$ 10 UIT. \\
Notaría y SUNARP & Lo habitual: los compradores, salvo que pacten mitad o que Richar asuma. Marquen el pacto en la minuta. \\
Planos / ingeniero & Quien encargue el levantamiento (en la práctica, los compradores). \\
\end{tabularx}}

\vspace{0.35em}
\noindent Estos cuatro ítems \textbf{no se ``legalizan'' con el voucher de efectivo de 2018}. Se pagan o se acreditan ahora, al formalizar.

\LugarFecha
\FirmasCompradores
\vspace{0.4em}
\noindent
\BloqueFirma{EL VENDEDOR}{\VendedorNombreCompleto}{\VendedorDNI}

\end{document}
